Trial-level scanpath descriptors are the workhorse of gaze-based classification \citep{selim2024review}. Recurrence quantification analysis summarizes re-inspection \citep{anderson2013recurrence}; transition matrices and their graph-theoretic reductions summarize the ordered structure of movement between regions \citep{smith2010eye, popelka2016scangraph}; entropy measures summarize the dispersion and predictability of looking; and mean fixation duration and saccade amplitude remain in nearly every feature set \citep{holmqvist2011eye}. When one of these descriptors discriminates a task \citep{greene2012reconsidering, borji2014defending, coutrot2018scanpath, shubi2025decoding}, an expertise level \citep{gegenfurtner2011expertise}, or a cognitive state, the result is usually reported as an accuracy. What is rarely reported is which properties of the scanpath the discrimination rests on: the set of locations, their order, how long each was held, or where on the stimulus they fell. Different descriptors encode these properties in different mixtures, and a descriptor's name is a poor guide to what it uses in a given dataset.

The consequences are concrete. In one of the public datasets analysed below, a classifier separates programmers who answered a comprehension question correctly from those who did not with an AUC of 0.72, using only the extent and shape of their fixations' spread. Most wrong answers came after one of the two programs, whose code has a different shape on the screen; when the answers are shuffled within each program the same classifier still averages 0.71, and among trials of the same program the spread carries no detectable signal. An accuracy cannot reveal this. Asking how the discrimination changes when a stated property of every scanpath is destroyed and the classifier is retrained, and comparing the answer with simple baselines, can.

Perturbing scanpaths to find out what a measure is sensitive to has precedent. \citet{dewhurst2012multimatch} validated MultiMatch with synthetic spatial-noise and order manipulations, and \citet{fahimi2021metrics} formalized this into a family of perturbations applied across a catalogue of scanpath similarity metrics to characterize their robustness. Shuffled-order baselines are standard in recurrence analysis \citep{anderson2013recurrence, anderson2015comparison}, static-versus-sequential comparisons appear in task decoding \citep{greene2012reconsidering}, and ordered-versus-shuffled gaze sequences have been used to dissociate cortical representations \citep{meghanathan2026temporal}. Outside eye tracking, retraining after removing information is a standard way to measure predictive importance \citep{hooker2019benchmark}, with formal inference for the resulting loss in predictiveness \citep{williamson2023general}. The scanpath studies ask how robust a \emph{similarity metric} is, or whether order matters for one task, and the importance methods operate on features rather than on scanpaths. Neither provides a general procedure for asking which properties of a scanpath a \emph{classification pipeline} depends on, with transformations whose invariants are stated, participant-aware evaluation, and uncertainty that has been checked.

This paper provides that procedure, the \emph{perturbation profile}, and makes four contributions. (1) Five transformations, each targeting one property of a scanpath (fixation order, the assignment of durations to fixations, the duration distribution, orientation and handedness, and fine spatial structure) while preserving a stated list of invariants, including the properties a transformation's name suggests it removes but does not. (2) A refit participant bootstrap for the change in participant-held-out AUC, whose coverage we check by simulation at the sample sizes analysed, together with exact invariance checks that catch implementation errors, a within-fold AUC that guards against evaluation bias, support diagnostics for descriptors defined on a bounded region, and a stated rule for which sensitivities are interpreted. (3) Demonstrations on thirteen problems from three public datasets (reading, program comprehension, and video and game viewing) in which the profile explains, with an exact result, why a descriptor improves when its own information is destroyed, helps identify a stimulus confound, and shows which sensitivities survive a change of evaluation metric, fold partition, Monte Carlo settings and learner. (4) An open Python package that profiles any descriptor with any scikit-learn learner. Section~\ref{sec:guide} turns the results into practical guidance.

The quantity measured throughout is the sensitivity of a refitted pipeline (descriptor, learner and participant-held-out evaluation together) to the destruction of a stated property. It is not an attribution of what a fitted classifier relied on, and it is not a measure of what participants did; the datasets below are demonstrations of that quantity rather than tests of hypotheses about viewing behaviour. The profile is meant for anyone who reports a scanpath-based classification and wants to say what the result rests on, and for anyone choosing between descriptors.
